\chapter{Research Objectives and Contributions}
\label{ch:contributions}

Our primary objective is to empirically evaluate the execution limits and scaling behavior of classical game-tree search heuristics implemented within managed, garbage-collected runtimes. Rather than seeking to compete with native C/C++ engines, we design Kronos as a controlled research platform to isolate and quantify the systems-level cost and benefit of individual search optimizations. This chapter details our core research questions, outlines the classical algorithms adopted from computer chess literature, and summarizes our main engineering contributions.

\section{Research Questions}
We formulate three core research questions:
\begin{enumerate}
    \item Can a chess search engine written in Node.js and JavaScript execute classical algorithms efficiently under V8's single-threaded and memory-managed runtime constraints?
    \item What is the relative search space reduction and playing strength contribution of individual heuristics (LMR, NMP, PVS, move ordering, and quiescence search) under controlled tournament conditions?
    \item What memory layout and object-allocation strategies are necessary and sufficient to prevent garbage-collection latency spikes during continuous execution?
\end{enumerate}

\section{Classical Techniques Adopted}
To isolate dynamic runtime effects, we adopt several standard search and evaluation heuristics from computer chess literature, claiming no algorithmic novelty for these components:
\begin{itemize}
    \item \textbf{Search Core:} Negamax formulation of Alpha-Beta pruning~\cite{knuth1975analysis} and Principal Variation Search (PVS)~\cite{marsland1982parallel}.
    \item \textbf{Search Reductions:} Null Move Pruning (NMP)~\cite{donninger1993null} and Late Move Reductions (LMR).
    \item \textbf{Move Ordering:} Most Valuable Victim -- Least Valuable Attacker (MVV-LVA), Killer Move cache, and History Heuristic~\cite{schaeffer1989history}.
    \item \textbf{Transposition Caching:} Zobrist Hashing~\cite{zobrist1970new} for state identification.
    \item \textbf{Positional Evaluation:} Piece-Square Tables (PSTs) with tapered interpolation.
\end{itemize}

\section{Engineering Contributions}
The contributions of this work lie in the runtime engineering, memory optimization, and deterministic benchmarking infrastructure designed to evaluate these heuristics:
\begin{itemize}
    \item \textbf{GC-Neutral Search Architecture:} An implementation that bounds V8 heap usage during tournament play (detailed in Chapter~\ref{ch:experimental_results}) by avoiding hot-path heap allocations via direct in-place board state mutations, flat TypedArray move buffers, and size-capped transposition table eviction.
    \item \textbf{Web Worker Concurrency Model:} An asynchronous coordination architecture separating the React user interface thread from concurrent background search threads (Web Workers) to prevent UI thread starvation.
    \item \textbf{Deterministic Benchmarking Pipeline:} A CLI-driven tournament suite executing paired, color-flipped matches using LCG-based opening book shuffles, ensuring identical game sequences across different host hardware configurations.
    \item \textbf{SPRT and Bradley-Terry Integration:} An automated verification framework utilizing Sequential Probability Ratio Testing (SPRT) bounds and Bradley-Terry rating regression to assess playing strength variations.
    \item \textbf{Incremental Checkpoint Scheduler:} A stateful scheduling architecture that logs tournament telemetry and matches incrementally to disk, enabling long-running evaluations to survive execution interrupts.
\end{itemize}


